\documentclass[11pt]{article}
\usepackage[margin=1in]{geometry}
\usepackage[T1]{fontenc}
\usepackage{lmodern}
\usepackage{hyperref}
\usepackage{booktabs}
\usepackage{longtable}
\usepackage{array}

\hypersetup{
  colorlinks=true,
  linkcolor=blue,
  urlcolor=blue,
  pdftitle={Rollout Statistics Module Audit},
  pdfauthor={Murphy}
}

\title{Rollout Statistics Module Audit}
\author{Murphy}
\date{2026-03-14 22:35 UTC}

\begin{document}
\maketitle
\sloppy

\section*{Scope}
This audit covers the dataset-rollout statistics path that produced the five-dataset
Qwen/Qwen3-1.7B bundle: \texttt{scripts/collect\_model\_stats.py}, the math / GPQA /
MMLU-Pro / LiveCodeBench adapters under \texttt{src/loop\_probe/adapters/},
\texttt{src/loop\_probe/prompt\_builder.py}, \texttt{src/loop\_probe/collector.py},
and \texttt{slurm/run\_collect\_model\_stats.sbatch}. I also asked Athena for an
independent audit of the same module and cross-checked the returned claims against the
local code and the generated summary bundle.

\section*{Executive Assessment}
Athena and the local read converge on the same conclusion: the current collector is
useful for descriptive loop-vs-length overlap analysis, but it is not yet a clean,
benchmark-faithful measurement stack for Qwen/Qwen3-1.7B. Three issues are definite
bugs or semantics bugs:

\begin{enumerate}
\item The math / GPQA / MMLU-Pro adapters leave Qwen3 in default thinking mode while
the rollout launch stays greedy (\texttt{temperature=0.0}).
\item The reported \texttt{max\_length\_hit} statistic is really an
``effective remaining context budget exhausted'' signal, not an unambiguous
``hit the model's native max length'' signal.
\item Prompt-too-long samples are counted differently on QA tasks versus
LiveCodeBench.
\end{enumerate}

A fourth high-confidence issue is the multiple-choice grader: GPQA and MMLU-Pro score
the \emph{last} standalone answer letter anywhere in the completion, which is brittle
when the model emits thought traces or extra commentary. The math verifier itself
looks sound.

\section*{Current Bundle Facts}
The regenerated five-dataset summary confirms that the active run used one common
generation policy: Qwen/Qwen3-1.7B, \texttt{temperature=0.0},
\texttt{max\_tokens=81920}, and \texttt{max\_model\_len=40960}. All five dataset rows
in \texttt{cross\_dataset\_rollout\_summary.json} report
\texttt{prompt\_too\_long = 0}, so the prompt-too-long inconsistency did not fire in
the shipped bundle; it is still a real latent bug for future runs.

\section*{Findings}
\begin{enumerate}
\item \textbf{Definite bug: Qwen3 default thinking-mode prompts are paired with greedy
decoding on the chat-template adapters.}

The math, GPQA, and MMLU-Pro prompt builders all call
\texttt{tokenizer.apply\_chat\_template(..., add\_generation\_prompt=True)} without
explicitly setting \texttt{enable\_thinking=False}
(\texttt{src/loop\_probe/prompt\_builder.py:4-10},
\texttt{src/loop\_probe/adapters/multiple\_choice\_gpqa.py:87-102},
\texttt{src/loop\_probe/adapters/multiple\_choice\_mmlupro.py:94-112}). The rollout
path then samples with \texttt{temperature=0.0}
(\texttt{scripts/collect\_model\_stats.py:450-458}), and the Slurm wrapper keeps that
as the default (\texttt{slurm/run\_collect\_model\_stats.sbatch:94-100}). The Qwen3
model card says reasoning/thinking is enabled by default and explicitly warns against
greedy decoding in thinking mode.\footnote{\url{https://huggingface.co/Qwen/Qwen3-1.7B}}

Why it matters: these three datasets can overstate loopiness or depress accuracy for
configuration reasons rather than task reasons. This is not a cosmetic prompt issue;
it changes the measurement contract itself.

Recommended fix: either set \texttt{enable\_thinking=False} for this collector family,
or switch the collector to a Qwen-recommended non-greedy decoding policy and document
it as intentional.

\item \textbf{Definite semantics bug: \texttt{max\_length\_hit} does not mean what the
report language suggests.}

The collector defines
\[
\texttt{effective\_max} = \min(\texttt{max\_tokens},\ \texttt{max\_model\_len} - \texttt{prompt\_len})
\]
in \texttt{scripts/collect\_model\_stats.py:152-153}, and marks a hit only when
\texttt{finish\_reason == "length"} and \texttt{token\_count == effective\_max}
(\texttt{scripts/collect\_model\_stats.py:479-490}). With the shipped defaults
\texttt{max\_tokens=81920} and \texttt{max\_model\_len=40960}
(\texttt{slurm/run\_collect\_model\_stats.sbatch:99-100}), the operative cap is
usually the configured remaining context budget, not the nominal \texttt{max\_tokens}.
vLLM's own docs define \texttt{max\_model\_len} as the model context length for
prompt plus output.\footnote{\url{https://docs.vllm.ai/en/latest/serving/engine_args.html}}

Why it matters: the current report wording ``hit max model length'' can be read as
``the model hit its native context limit'' even though the implementation is only
detecting exhaustion of the configured local cap. This is especially risky because the
Qwen3-1.7B model card advertises a 32{,}768-token context window, while the collector
is configured to 40{,}960.

Recommended fix: rename the metric to something like
\texttt{effective\_context\_cap\_hit}, log both prompt length and configured cap, and
either justify the 40{,}960 override or reduce it to a validated value.

\item \textbf{Definite bug: prompt-too-long rows are scored inconsistently across
tasks.}

When \texttt{effective\_max < 1}, QA tasks increment
\texttt{num\_prompt\_too\_long} and then skip grading entirely
(\texttt{scripts/collect\_model\_stats.py:426-447}). Since
\texttt{success\_fraction = num\_correct / num\_graded}
(\texttt{src/loop\_probe/collector.py:115-116}), those rows vanish from QA accuracy.
LiveCodeBench does something different: it still writes an
\texttt{LcbSampleRecord} with empty code and passes that record into native grading
(\texttt{scripts/collect\_model\_stats.py:433-445},
\texttt{src/loop\_probe/adapters/livecodebench\_codegen.py:253-300}), which
effectively counts the row as wrong.

Why it matters: the same failure mode inflates QA success rates but not
LiveCodeBench success rates. Even though this branch's current bundle had zero such
rows, the cross-task comparison contract is inconsistent in code.

Recommended fix: either exclude prompt-too-long rows from grading on \emph{all}
tasks, or count them as wrong on \emph{all} tasks. The branch should pick one policy
and serialize it in the report metadata.

\item \textbf{High-confidence grading bug path: GPQA and MMLU-Pro parse the last
standalone option letter anywhere in the response.}

Both graders regex the completion for option letters and score \texttt{matches[-1]}
(\texttt{src/loop\_probe/adapters/multiple\_choice\_gpqa.py:105-109},
\texttt{src/loop\_probe/adapters/multiple\_choice\_mmlupro.py:115-135}). The prompts
only ask for ``the single best answer letter'' and do not force a structured terminal
field (\texttt{src/loop\_probe/adapters/multiple\_choice\_gpqa.py:93-102},
\texttt{src/loop\_probe/adapters/multiple\_choice\_mmlupro.py:103-112}).

Why it matters: if the model emits reasoning, re-lists options, or leaves a thought
trace before the final answer, the scorer can assign the wrong label even when the
model ``meant'' the correct option. The earlier thinking-mode bug increases the
chance of exactly this failure mode.

Recommended fix: require a structured final field (for example a terminal
\texttt{Final answer: X} or JSON answer slot) and parse only that field. One caveat:
the current collector does not checkpoint raw MCQ completions, so this bug path is
proven by the parser logic but cannot be quantified retroactively on the shipped
bundle.

\item \textbf{Risk / assumption: the LiveCodeBench path is intentionally non-canonical
but is easy to read as canonical.}

The adapter maps every \texttt{qwen/*} model to
\texttt{LMStyle.CodeQwenInstruct}
(\texttt{src/loop\_probe/adapters/livecodebench\_codegen.py:181-193}), which is a
reasonable guess but still a local assumption for Qwen3-1.7B. The collector also
hard-codes one sample per prompt (\texttt{n=1}) and the current launch policy keeps
\texttt{temperature=0.0}
(\texttt{scripts/collect\_model\_stats.py:450-458},
\texttt{slurm/run\_collect\_model\_stats.sbatch:94-100}). Upstream LiveCodeBench code
generation instructions instead use \texttt{n=10}, \texttt{temperature=0.2}, and the
parser defaults to \texttt{num\_process\_evaluate=12} and \texttt{timeout=6}; the
local adapter defaults to \texttt{1} and \texttt{5}
(\texttt{src/loop\_probe/adapters/livecodebench\_codegen.py:100-118}).\footnote{\url{https://github.com/LiveCodeBench/LiveCodeBench/blob/main/README.md}}
\footnote{\url{https://github.com/LiveCodeBench/LiveCodeBench/blob/main/lcb_runner/runner/parser.py}}

Why it matters: the current LiveCodeBench number is perfectly usable as deterministic
task-local telemetry, but it is not benchmark-canonical. The report should say so
explicitly.

Recommended fix: either align generation and evaluation with the upstream benchmark
contract, or label the output as a local single-rollout deterministic variant.
\end{enumerate}

\section*{What Looks Correct}
\begin{itemize}
\item \textbf{Math verification is in good shape.} The math prompt asks for a boxed
final answer (\texttt{src/loop\_probe/prompt\_builder.py:4-10}), and
\texttt{math\_freeform.py} uses \texttt{math\_verify} with boxed-priority extraction
plus expression fallback (\texttt{src/loop\_probe/adapters/math\_freeform.py:24-91}).
The verifier call order is also correct: \texttt{verify(parsed\_gold, parsed\_pred)}
at \texttt{src/loop\_probe/adapters/math\_freeform.py:83-85}.
\item \textbf{The overlap counters themselves are straightforward.} Once a completion
is generated, the bookkeeping for loop counts, length sums, and overlap counts in
\texttt{scripts/collect\_model\_stats.py:492-531} and
\texttt{src/loop\_probe/collector.py:114-156} is internally coherent. The dominant
risks are in the prompt/decode/grading contract around those counters, not in the
counter arithmetic.
\end{itemize}

\section*{Recommended Repair Order}
\begin{enumerate}
\item Fix the Qwen3 prompt/decode contract first. This is the highest-impact issue for
all chat-template datasets.
\item Rename and re-document the ``max length hit'' metric before anyone treats it as
a literal native-model ceiling statistic.
\item Unify prompt-too-long handling across QA and LiveCodeBench.
\item Replace MCQ letter scraping with a structured terminal-answer parser and add raw
completion checkpointing for auditability.
\item Decide whether LiveCodeBench is meant to be canonical benchmark evaluation or a
local deterministic telemetry variant, then encode that decision in the report text
and runtime defaults.
\end{enumerate}

\section*{Bottom Line}
The current five-dataset bundle is still informative for one descriptive claim:
looped rollouts and effective context-cap terminations overlap strongly under this
runtime. What it does \emph{not} yet provide is a clean apples-to-apples statement of
Qwen3-1.7B task accuracy or canonical benchmark performance, because the prompt,
decode, and grading contracts are not yet consistent enough for that stronger claim.

\end{document}
